% TOP_PROBLEM: {"id": 4700028, "title": "Around Kouchnirenko's conjecture I", "category_id": 4, "category": {"id": 4, "name": "algebra", "display_name": "Algebra", "description": "Group theory, ring theory, field theory, and algebraic structures.", "slug": "algebra", "order_index": 4, "created_at": "2026-07-31T15:26:25.671Z"}, "status": "partially_solved", "problem_number": "AMR-046-0028", "set_id": 13}
% FIRST_POSTED: 2026-10-02
% ENTRY_KIND: statement_only
% UnsolvedMath ID 4700028; source catalogue code AMR-046-0028
% !TeX program = lualatex
% Definition and sources only; this file is not an LLM solution attempt.
\documentclass[11pt,a4paper]{article}
\usepackage[margin=25mm]{geometry}
\usepackage{fontspec}
\setmainfont{lmroman10-regular.otf}[BoldFont=lmroman10-bold.otf,
 ItalicFont=lmroman10-italic.otf,BoldItalicFont=lmroman10-bolditalic.otf]
\usepackage{amsmath,amssymb,mathtools}
\usepackage{unicode-math}
\setmathfont{latinmodern-math.otf}
\AtBeginDocument{\let\setminus\smallsetminus}
\usepackage{xurl,hyperref}
\hypersetup{colorlinks=true,urlcolor=blue}
\setcounter{secnumdepth}{0}
\setlength{\parindent}{0pt}
\setlength{\parskip}{5pt}
\setlength{\emergencystretch}{3em}
\begin{document}
\section*{Around Kouchnirenko's conjecture I}\label{4700028}
{\small UnsolvedMath ID 4700028 · AMR-046-0028 · Algebra · AMR Open Problem Lists}
\subsection{Definitions and mathematical statement}
Fix positive integers $m_1,m_2$. Consider two real polynomials
$f_i\in\mathbb R[x,y]$, each with at most $m_i$ nonzero monomial terms.
The exponents are nonnegative integers, with no degree bound. A common
zero $(u,v)$ is simple, or nondegenerate, when
\[
 f_1(u,v)=f_2(u,v)=0,\qquad
 \det\begin{pmatrix}(f_1)_x&(f_1)_y\\(f_2)_x&(f_2)_y\end{pmatrix}(u,v)\ne0.
\]
Let $N_+(m_1,m_2)$ be the supremum of the number of such zeros in the
strictly positive quadrant $(0,\infty)^2$ over all allowed polynomial pairs.
The problem asks for a useful, ideally sharp, upper bound for
$N_+(m_1,m_2)$ depending only on the term counts, and examples showing how
close the bound is to optimal.

Zeros on either coordinate axis are excluded. Multiplicities are not
counted, and positive-dimensional common-zero sets do not contribute
nondegenerate zeros. A common monomial factor may be removed without
changing the positive solutions, illustrating why total degree is not
the relevant measure of complexity.

The naive product $(m_1-1)(m_2-1)$ is not a valid general upper bound:
the source discusses counterexamples and the sharp five-solution result
for two trinomials. General fewnomial estimates give finite bounds, but
the question seeks stronger estimates reflecting the separate term
counts. A Bézout estimate that grows with unrestricted degrees does not
provide the requested uniform control.
\subsection{Short English statement}
Find sharp or substantially improved bounds for positive simple zeros of two polynomials using only their numbers of terms.
\subsection{Sources}
\begin{itemize}
\item[S1] ULAM.AI and the UnsolvedMath Contributors, supplied problems.json, record 4700028 (AMR-046-0028), AMR Open Problem Lists. Catalogue identity and source transcription; CC BY 4.0.
\url{https://huggingface.co/datasets/ulamai/UnsolvedMath}.
\item[S2] Gasull - Some open problems in low dimensional dynamical systems (2020), Problem environment 28: Around Kouchnirenko's conjecture I. Original source indexed by AMR-046-0028.
\url{https://arxiv.org/abs/2012.02524}.
\end{itemize}
\end{document}
